\section{Introduction}\label{sec:intro}

A large language model asked to write a unit test for a particular function often fails to
produce compilable code. It is tempting to blame the model, and we read it the same way in an
earlier analysis of \texttt{gpt-4o-mini} on 120 medium-complexity functions. This paper
reports on the investigation that led us to that conclusion and explains why the answer is
mostly about a benchmarking artifact rather than the model itself.

We found two issues, and together they compose a problem.

\paragraph{Benchmarks contain targets that no generator can test.}
Sampling functions by cyclomatic complexity cuts off the tail, but nothing else. If the
function is not reachable from an external test, neither can the generator hope to test it,
and asking it to do so will label it as failed. We found that in the ten projects studied,
1{,}848 functions fell in the desired complexity range of $5 \le CC \le 10$. Less than a
third of them (474) were simultaneously public and had a clear name, suggesting that a
generator that failed on the remaining 74\,\% would have been asked to do what no tool was
capable of doing. Likewise, attribution of a method to the wrong declaring class renders a
test suite invalid, and there is ample opportunity for that to happen: a brace counter that
does not decrement on a method exit will assign the method to whichever nested type was last
opened.

\paragraph{The usual metrics can be satisfied with no testing at all.}
Compilation success and coverage are both metrics that can be reported by a suite that never
tests anything. Reflection turns every identifier into a string, so any call to a method that
does not exist compiles cleanly with a zero return code; and a generated suite in our own
corpus contained a single assertion-free call to the wrong function, which passed. Worse, the
suite's metrics were all good on a superficial inspection. The issue was not one of noisy
metrics; instead, the suite passed because it asserted nothing at all. We found six such
cases, and in four of them, the defect was ours when repairing the other two
(Section~\ref{sec:threats}).

\subsection*{This paper}

We build a new benchmark satisfying five testability criteria established in advance, sample
60 Java and 60 Python functions from the resulting set, and evaluate each test suite at four
nested levels. We rely exclusively on mutation score within the line range as a metric. All
other claims are vacuous for the suite that finds nothing because it detects no faults. Two
variants of the prompt differ in the information given about the target, allowing us to report
on the contribution of the target specification, while three established tools (EvoSuite,
Randoop, Pynguin) act as reference points for comparisons. The details of sampling, tool
versions, budget, and hypotheses were fixed in a public repository before measurements began.

Three findings stand out.

\begin{enumerate}
  \item \textbf{Specification improves reach, not precision.}
    Adding a fully qualified target name, its signature, and a working constructor call to
    the prompt increases the proportion of suites reaching the target in Python from 29/60 to
    42/60 ($p = 0.003$, rank-biserial $+0.61$). The number detecting a fault rises from 16/60
    to 22/60, but this is not significant ($p = 0.33$). The prompt gets the test to the right
    place; it does not make it check anything. A study that used coverage as its sole metric
    would have announced the change as an improvement.
  \item \textbf{LLMs and search-based tools have their own merits.}
    The prompt conditions that add the target specification beat Pynguin across the board in
    Python. With Pynguin detecting a fault on only 2/60, compared to 16/60 and 22/60 for the
    two prompt conditions, they are better in every pairing at $p \le 0.003$, and this applies
    to both of the configurations we tried for Pynguin. The trend is reversed in Java, where
    both search-based methods detect faults on more targets than the model --- significantly
    so only for EvoSuite, and only under the stricter gate. These findings do not support one
    particular approach.
  \item \textbf{Most of the attrition is structural.}
    The issue was never with writing assertions. It was with getting an instance of the
    subject under test, getting any assertions to pass, and generally navigating the maze of
    object creation that the search-based literature has already acknowledged for years. That
    is the single challenge that the model was unable to bypass.
\end{enumerate}

% MUC MOI — van cua AI, tac gia can viet lai bang giong cua minh truoc khi nop (RBL-5b).
\subsection*{Contributions}

\begin{itemize}
  \item \textbf{A quantified testability funnel.} Only 26\,\% of the medium-complexity
    functions in ten real projects are addressable by an external test at all. Any benchmark
    that samples on complexity alone charges its tools for the other three quarters. We give
    the five criteria and the per-stage counts so the funnel can be reproduced elsewhere.
  \item \textbf{A four-tier measurement that cannot be satisfied vacuously.} We show, with
    cases drawn from this corpus, that compilation success and a green-test criterion are both
    satisfiable by suites that verify nothing, and we rest every claim on the one tier that is
    not.
  \item \textbf{Evidence that coverage and fault detection can move apart.} A single, plausible
    prompt change raises the proportion of suites reaching the target significantly while
    leaving fault detection unchanged --- far enough apart to invert the conclusion a
    coverage-only study would draw.
  \item \textbf{A cross-language warning.} Measured with one harness, the model beats the
    Python baseline and trails both Java baselines. A study confined to either language would
    have reported the opposite of the other.
  \item \textbf{A catalogue of six ways a measurement fails silently}, four of them introduced
    by us while repairing the other two, together with the defensive checks we now apply. One
    of the six, had we trusted it, would have supported a statistically significant claim in
    the wrong direction.
\end{itemize}

We report the full sample size at all stages of the benchmarking funnel, and we discuss the
two configurations of Pynguin at length. We describe an anomaly in the Java implementation of
a green-test gate --- the criterion requiring at least one test to pass against the unmodified
source --- that we found to contradict its own pre-registered specification. Where a repair
procedure favors us, we report the effect size and the outcome in both configurations, so that
nothing in this paper should depend on the reader's preference between the two.
